\section{Discussion}
\label{sec:discussion}

\subsection{The Spike-vs-Analog Gap}

The most striking architectural finding is the gap between spike-based and analog readout, documented in Section~\ref{sec:exp_mnist}. The same trained network produces chance accuracy via a spike-based motor readout and high accuracy via an analog readout over the same morphon potentials. The information is present in the potentials --- the analog readout proves it --- but the spike pipeline destroys it through propagation delays, leaky integration, and multi-hop conversion.

This is not a bug in MI; it is a property of any spiking architecture in which the readout pathway sums potentials through several intermediate firing decisions. Biology solves this by having dedicated analog readout pathways for tasks that require it (e.g.\ Purkinje cells in the cerebellum perform analog integration of climbing fiber and parallel fiber inputs to produce a continuous motor signal; \cite{ito2006cerebellar}). MI's analog readout is the in-silico analogue.

The spike pipeline still has a job: it produces the diverse, sparse, energy-efficient hidden representations that the analog readout exploits. The two pathways are complementary, not competing. The MI thesis is not ``replace spikes with analog'' but ``spike-based developmental dynamics produce features that analog readouts can use.''

\subsection{Self-Healing as a Diagnostic}

The v3.0.0 codebase showed what appeared to be beneficial self-healing: a 30\% intact network recovered to 48\% after 30\% ablation, an 18pp gain over its pre-damage state. This result is real and reproducible in that codebase. In v4.1.0 it disappears: the intact network reaches 52\% and post-damage recovery falls to 28.5\%.

The mechanism is now clear. In v3.0.0, rate-based iSTDP failed to suppress hub morphons during training. Hubs entrenched while Endoquilibrium's Mature stage gate froze plasticity, leaving the system stuck at 30\%. Damage removed hubs at random; Endoquilibrium re-entered the Differentiating stage in response to the performance drop, providing the high-plasticity window the original training never sustained; regrowth filled the gap with well-specialized morphons. The +18pp came entirely from two sources: hub removal and plasticity restoration --- exactly what spike-timing iSTDP provides automatically.

This reframes the v3.0.0 self-healing result as a diagnostic rather than a feature. It told us: (1) hubs are the bottleneck; (2) the substrate is capable of 48\% if given the right plasticity regime. Spike-timing iSTDP is the engineered solution to both. The v4.1.0 intact ceiling (52\%) is not coincidentally above the v3.0.0 post-recovery ceiling (48\%); it is the same substrate, finally given the conditions it needed.

The general principle is: \emph{when a developmental system benefits from damage, the damage is revealing a regulatory failure, not demonstrating a useful property}. The right response is to find and fix the failure, not to incorporate damage into the training protocol.

\subsection{Limitations}

We are honest about the gaps:

\begin{itemize}
  \item \textbf{MNIST accuracy is well below state of the art.} 52\% intact (v4.1.0 standard profile) is a large improvement over 30\% (v3.0.0) but still far from Diehl \& Cook's $\sim$95\% STDP baseline \cite{diehl2015unsupervised} and from supervised SNNs trained with surrogate gradients. Hub dominance is broken; the remaining gap is in feature quality and readout depth. Parameter tuning (\texttt{istdp\_rate}, $\alpha$, inhibitory weight initialization) and a deeper readout are the next targets.
  \item \textbf{No language tasks yet.} The system can encode and retain information across sequential inputs (temporal memory exists in residual potentials), but compositional reasoning is not yet demonstrated. Composition requires nonlinear hidden features formed under DFA-driven supervised feedback, which the current spike pipeline does not deliver.
  \item \textbf{The premature Mature fix is a patch, not a principled solution.} The history gate at 2000 ticks prevents the failure but is essentially ``disable Mature stage detection.'' The principled fix is RPE-based regulation via a Limbic Circuit module that does not yet exist.
  \item \textbf{Performance.} The v4.1.0 implementation processes $\sim$23 images/second with $\sim$74K synapses (130~s for 3000 images), a 4.4$\times$ improvement over v3.0.0 via Pulse Kernel Lite. The medium path (eligibility traces, weight updates) is now the dominant cost; sparse eligibility updates and parallel medium-path learning are the next target.
\end{itemize}

\subsection{The Path to Language}

The planned trajectory beyond this paper:
\begin{enumerate}
  \item \textbf{Local inhibitory competition with iSTDP \cite{vogels2011inhibitory}.} \textit{Implemented in v4.1.0.} Hub dominance broken (1 labeled class $\rightarrow$ 12, accuracy 8\% $\rightarrow$ 44.5\%). Spike-timing iSTDP per delivered spike; 4.4$\times$ wall-clock speedup via Pulse Kernel Lite hot arrays. Next: parameter tuning toward $\geq$85\%.
  \item \textbf{Limbic Circuit + RPE-based regulation.} Replace the history-gate patch with reward prediction error feeding a Motivational Drive module. This naturally handles both RL and classification.
  \item \textbf{Temporal sequence processing.} Recurrent connectivity within associative populations + working-memory feedback into sensory ports. Spec exists; implementation gated on MNIST $\geq$ 85\%.
  \item \textbf{Tokenization and embedding.} Once the temporal pipeline works, add a learned tokenizer and vocabulary embedding. This is the bridge to actual NLP tasks.
\end{enumerate}

The structural decisions for each step are documented in the project's specification corpus. The question this paper does \emph{not} attempt to answer is whether MI will ever match Transformer performance on language tasks. We argue only that developmental dynamics are a principled and underexplored axis of neural architecture design, and that the failure modes we encountered have biological analogues whose biology-derived fixes work.
